\section{\texorpdfstring{$\mathbb{R}$ e la cardinalit\`a del continuo}{R e la cardinalit\`a del continuo}}

In questa sezione daremo una definizione di $\RR$ come insieme ordinato. Estenderemo, poi, la definizione ad includere le operazioni di campo, ma senza svolgere le verifiche.

\begin{definition}
	Sia $(A,<)$ un ordine totale, allora:
	\begin{itemize}
		\item $m \in A$ \`e un \vocab{maggiorante} di $B \subseteq A$ se $\forall {x \in B} \; x \leq m$
		\item $B \subseteq A$ \`e \vocab{superiormente limitato} se ha un maggiorante
		\item $s\in A$ \`e \vocab{l'estremo superiore} di $B$ -- denotato $\drkblu \sup B$ -- se $s$ \`e il minimo dei maggioranti di $B$.
	\end{itemize}
\end{definition}

\begin{note}
	Non sempre gli estremi superiori esistono, per\`o, se $B$ ha un estremo superiore, questo \`e unico.
\end{note}

\begin{definition}
	Un ordine totale $(A,<)$ \`e \vocab{completo} se ogni $B \subseteq A$ non vuoto e superiormente limitato ha un estremo superiore $\sup B \in A$.
\end{definition}

\begin{exercise}
	\label{q_noncompleto}
	Dimostra, usando solo le propriet\`a di $\QQ$, che l'insieme $\{x
\in \QQ | x^2 < 2\}$ non ha estremo superiore in $\QQ$.

{\drkred Attenzione!} Dire che il sup \`e $\sqrt 2$ ma $\sqrt 2\not\in\QQ$ non
vale. Tutto quello che hai \`e~$\QQ$.
\end{exercise}

In conseguenza dell'esercizio, possiamo dire che $\QQ$ non \`e completo. Costruiamo ora un ordine completo $(\RR,<)$ che contiene una copia isomorfa di $\QQ$ come sottoinsieme denso.

\begin{definition}
	Sia $(A,<)$ un ordine totale. $B \subseteq A$ \`e un
\vocab{segmento iniziale} di $A$ se
\[\forall {x \in B} \; \forall {y \in
A} \; y < x \rightarrow y \in B\]
{\drkgrn Ossia $B$ \`e un segmento iniziale di $A$ se, ogniqualvolta $B$
contiene un elemento, $B$ contiene altres\`i tutti gli elementi minori di
questo.}

Un segmento iniziale $B$ di $A$ si dice \vocab{proprio} se $B \ne A$.
\end{definition}

\begin{example}[Segmento iniziale principale]
	Dato $(A,<)$ ordine totale, $A$ stesso e $\emptyset$ sono segmenti iniziali di $A$. Dato $x \in A$, l'insieme:
	\[ {\drkblu A_x}  \;\;\Mydef\;\; \{y \in A | y < x\}
		\]
	\`e un segmento iniziale proprio di di $A$, detto \vocab{segmento inziale principale} determinato da $x$.
\end{example}

\begin{note}
	Useremo nuovamente il concetto di segmento iniziale studiando gli ordinali. Il prossimo concetto, quello di sezione di Dedekind, invece, ci serve unicamente per definire $\RR$.
\end{note}

\begin{definition}[Sezioni di Dedekind]
	Una \vocab{sezione} sull'insieme totalmente ordinato $(A,<)$ \`e
un segmento iniziale {\drkred proprio}
	e {\drkred non vuoto} di $A$ che {\drkred non ha un massimo
elemento.}
{\drkgrn Ossia $B$ segmento iniziale di $A$ \`e una sezione se $B \ne A$,
$B \ne \emptyset$ e $\forall {x \in B} \; \exists {y \in B} \; x < y$.}
\end{definition}

\begin{definition}[$\RR$ come insieme ordinato]
	\[ {\drkred \RR} \;\;\Mydef\;\; \{\,x \in \ps(\QQ)\, | \,
\text{$x$ \`e una sezione su $\QQ$}\,\}
		\]
	Dati $x,y\in\RR$
	\[ {\drkred x \leq y} \;\;\Mydef\;\; x \subseteq y
		\]
\end{definition}

\begin{proposition}[$\RR$ \`e completo]\label{prop-r-comp}
	$(\RR,<)$ \`e un ordine totale completo.
\end{proposition}

Prima della dimostrazione, isoliamo un semplice lemma.

\begin{lemma}[L'unione di segmenti iniziali \`e un segmento iniziale]
	Sia $(A,<)$ un ordine totale e $X$ un insieme di segmenti iniziali di $A$. Allora $\bigcup X$ \`e un segmento iniziale di $A$.
\end{lemma}

\begin{proof}
	Sia $\alpha \in \bigcup X$ e sia $\beta \in A$, con $\beta < \alpha$, vogliamo verificare che $\beta \in \bigcup X$.
	Siccome $\alpha \in \bigcup X$, esiste $x\in X$ con $\alpha \in
x$. Siccome $x$ \`e segmento iniziale di~$A$, $\beta \in x \subseteq \bigcup X$.
\end{proof}

\begin{proof}[Dimostrazione di~\ref{prop-r-comp}]
	Abbiamo un ordine parziale perch\'e il contenimento ($\subseteq$)
\`e un ordine parziale. Dobbiamo verificarne la totalit\`a.
	Dati $x,y \in \RR$, supponiamo per assurdo che non valga n\'e $x \subseteq y$ n\'e $y \subseteq x$, allora esistono $a,b \in \QQ$ tali che $a \in x \setminus y$ e $b \in y \setminus x$.
	Ora, o $a < b$, o $b < a$.
	Nel primo caso, poich\'e $b \in y$, si avrebbe $a \in y
\;\lightning$. Analogamente, nel secondo caso, $a\in x \Rightarrow b \in x \; \lightning$.

	Verifichiamo ora la completezza. Sia $A \subseteq \RR$ non vuoto e
superiormente limitato da $m \in \RR$. Dimostriamo che $\sup A = \bigcup A \in \RR$.

{\drkblu $\bigcup A \in \RR$.}	Per il lemma precedente $\bigcup A$ \`e un segmento iniziale di
$\QQ$. Siccome $A$ non \`e vuoto, allora $\bigcup A \ne \emptyset$. Poich\'e $m$ \`e un maggiorante di $A$,
	si ha $\bigcup A \subseteq m$, per cui $\bigcup A \ne \RR$. Se
$\bigcup A$ avesse un massimo~$M$, allora esisterebbe $a\in A$ con $M\in
a$, per cui $M$ sarebbe il massimo anche di~$a\;\lightning$. In definitiva
$\bigcup A$ \`e una sezione di $\QQ$, quindi $\bigcup A \in \RR$.

{\drkblu $\bigcup A$ \`e un maggiorante di $A$.} Infatti, se $x \in A$,
allora $x \subseteq \bigcup A$, cio\`e $x\le\bigcup A$.

{\drkblu $\bigcup A$ \`e il minimo maggiorante.}
	Se $m$ \`e un maggiorante di $A$, allora per definizione $\forall
{x \in A} \; x \subseteq m$, ma ci\`o significa che $\bigcup A \subseteq
m$, ovvero $\bigcup A \le m$.
\end{proof}

\begin{remark}
	La funzione $\iota \colon \QQ \to \RR \quad a \mapsto \QQ_a
\drkgrn = \{x \in \QQ | x < a\}$
	immerge $\QQ$ in $\RR$ in maniera strettamente crescente e densa
-- ossia $\iota[\QQ]$ \`e denso in $\RR$.
\end{remark}

\begin{proof}
	Dati $a,b \in \QQ$ con $a < b$, abbiamo $\QQ_a \subsetneq \QQ_b$,
per cui $\iota$ \`e strettamente crescente. Dati $x,y \in \RR$ con $x <
y$, per definizione $x \subsetneq y$, quindi c'\`e un $a \in y \setminus
x$. Siccome $y$ non ha massimo, c'\`e un $b\in y$ con $a<b$. Ora,
$x\subseteq \QQ_a \subsetneq \QQ_b \subsetneq y$, per cui $x<\iota(b)<y$.
\end{proof}

\begin{notation}[Abuso di]
	Siccome le immersioni:
	\[ \omega \hookrightarrow \ZZ \hookrightarrow \QQ \hookrightarrow \RR
		\]
	sono tutte iniettive e crescenti, quando non c'\`e pericolo di confusione, possiamo abusare della notazione immaginando che 
	queste siano vere e proprie inclusioni:
	\[ \omega \;{\drkred\subseteq}\; \ZZ \;{\drkred\subseteq} \;\QQ
\;{\drkred\subseteq}\; \RR
		\]
	In realt\`a non \`e vero: per esempio \`e $\iota[\QQ]$, non $\QQ$, a essere sottoinsieme di $\RR$, ma $\iota[\QQ]$ \`e in corrispondenza biunivoca, in maniera 
	canonica, con $\QQ$, e questa corrispondenza preserva tutta la
struttura rilevante: l'ordine, come abbiamo verificato, ma anche le operazioni di campo.
\end{notation}

\begin{corollary}
	$\aleph_0 < |\RR|$.
\end{corollary}

\begin{proof}
	Abbiamo visto che $|\QQ| \leq |\RR|$, inoltre $\QQ$ \`e denso in $\RR$, pertanto 
	$\RR$ \`e denso in s\'e. Si vede facilmente che $\RR$ non ha
massimo n\'e minimo, quindi, se $\RR$ fosse numerabile, sarebbe isomorfo a
$\QQ$, per il \hyperref[isoCantor]{teorema di isomorfismo di Cantor}.
	D'altro canto $\RR$ \`e completo e $\QQ$ no per
l'esercizio~\ref{q_noncompleto}.
\end{proof}

\subsection{Caratterizzazione di $(\RR, <)$.}
Abbiamo stabilito che $(\RR,<)$ \`e un ordine totale, completo e senza estremi con un sottoinsieme, $\QQ$, denso e numerabile.
Queste propriet\`a, a loro volta, caratterizzano l'insieme ordinato $(\RR,<)$ a meno di isomorfismi.

\begin{proposition}[Caratterizzazione di $(\RR,<)$]\label{pro-car-rrr}
	Sia $(A,<)$ un ordine totale, allora $(A,<) \sim (\RR,<)$ se e
solo se valgono simultaneamente:
	\begin{enumerate}[1.]
		\item $(A,<)$ \`e completo
		\item $(A,<)$ \`e senza estremi
		\item esiste $B \subseteq A$ numerabile e denso in $A$.
	\end{enumerate} 
\end{proposition}

\begin{proof}
	Denotiamo con $\widetilde{A}$ l'insieme delle sezioni {\drkred su
$B$}.
	Osserviamo che $(A,\leq) \sim (\widetilde{A},\subseteq)$.
	L'isomorfismo \`e infatti dato da:
	\[ f \colon A \rightarrow \widetilde{A} \qquad a \mapsto B_a = \{x \in B | x < a\}
		\]
	la cui inversa \`e:
	\[ g \colon \widetilde{A} \rightarrow A \qquad Y \mapsto \sup Y
		\]

	Verifiche: dobbiamo mostrare che $g \circ f = \id_A$, $f \circ g =
\id_{\widetilde{A}}$ e $f$ \`e crescente.
	
{\drkblu $g \circ f = \id_A$} Dato $a \in A$, sia $b = g(f(a)) = \sup B_a
= \sup\{x \in B | x < a\}$. Dobbiamo dimostrare $a=b$. Per costruzione $a$
\`e un maggiorante di~$B_a$. Supponiamo che non sia il minimo, allora
			preso $c < a$ maggiorante di~$B_a$, essendo $B$
denso, esiste $c'\in B$ con $c < c' < a$, quindi $c' \in B_a\;\lightning$

{\drkblu $f \circ g = \id_{\widetilde{A}}$} Data una sezione $X$ di $B$,
sia $Y = f(g(X)) = B_{\sup X}$. Mostriamo che $X = Y$. Siccome $X$ non ha
massimo, $\forall {x \in X} \; x < \sup X$, 
			per cui $x \in B_{\sup X}=Y$, quindi $X \subseteq
Y$. Viceversa, dato $y\in Y$, ossia $y < \sup X$, esiste $y' \in X$ tale
che $y < y'$, per cui $y \in X$. Ne segue che  per cui $Y \subseteq X$.

{\drkblu $f$ \`e crescente.} Dati $a,a' \in A$, se $a<a'$, allora $B_a
\subsetneq B_{a'}$ perch\'e, per la densit\`a di~$B$, esiste $b\in B$ con
$a<b<a'$.

	Ora, per il \hyperref[isoCantor]{teorema di isomorfismo di
Cantor}, $B$, con l'ordine indotto da~$(A,<)$, \`e isomorfo a~$\QQ$,
quindi le sezioni di $B$ sono isomorfe a quelle di~$\QQ$, per cui $(A,<)
\sim (\widetilde{A},\subsetneq) \sim (\RR, <)$.
\end{proof}

Per completezza, definiamo la struttura di campo di $\RR$. Non verificheremo le propriet\`a, e neanche la correttezza di queste definizioni.

\begin{definition}[Campo ordinato]
	$(F,0,1,+,\cdot,\leq)$ \`e un \vocab{campo ordinato} se:
	\begin{itemize}
		\item $(F,0,1,+,\cdot)$ \`e un campo
		\item $(F,<)$ \`e un'ordine totale
		\item $\forall {x,y,z \in F} \; x < y \rightarrow x + z < y + z$
		\item $\forall {x,y \in F} \; (0 < x \land 0 < y) \rightarrow 0 < x \cdot y$
	\end{itemize}
\end{definition}

\begin{definition}
	Dati $x,y \in \RR$
	\[ {\drkblu x + y} \;\;\Mydef\;\; \{a + b \in \QQ | a \in x \land b \in y\}
		\]
	Dati $x,y \in \RR$ con $x > 0$ e $y > 0$
	\[ {\drkblu x \cdot y} \;\;\Mydef\;\; \{q \in \QQ | q \leq 0\} \cup \{a\cdot b \in \QQ| a \in x  \land b \in y \land a > 0 \land b > 0\}
		\]
Definiamo quindi $\drkblu -x$ tramite l'inverso additivo ed il prodotto
nei casi in cui uno dei fattori \`e $<0$ tramite la regola dei segni.
\end{definition}


\begin{theorem}[Unicit\`a di $(\RR,0,1,+,\cdot,\leq)$]
	$\RR$ dotato delle operazioni definite, \`e l'unico campo ordinato completo a meno di isomorfismo.
\end{theorem}

La dimostrazione di questo teorema, talvolta, si vede nei corsi di analisi
1, noi non la studieremo, Per chi fosse interessato, vedi per esempio
\textit{Hrbacek Jech, Introduction to Set Theory}, capitolo~10.

\subsection{$|\RR| = 2^{\aleph_0}$}
Torniamo ad una questione pi\`u strettamente insiemistica.

\begin{theorem}\label{crdr}
	$|\RR| = 2^{\aleph_0}$
\end{theorem}

Questo teorema ci dice, in un modo ancora diverso, che $\RR$ \`e pi\`u che
numerabile -- poich\'e $\aleph_0 < 2^{\aleph_0}$ -- ma, in pi\`u, caratterizza
anche esattamente la cardinalit\`a di $\RR$.

{\drkgrn Prima della dimostrazione formale, vediamo intuitivamente perch\'e il risultato \`e vero. Per definizione $\RR \subseteq \ps(\QQ)$, quindi la disuguaglianza da dimostrare \`e $2^{\aleph_0} \leq |\RR|$.
Esibiamo una funzione iniettiva $\ps(\omega) \to \RR$:
\[ f(S) = \underbrace{0.a_0^Sa_1^Sa_2^Sa_3^S\ldots}_{\text{in notazione
decimale}} \qquad \text{con} \; a_i^s = \begin{cases}
	0 &\text{\text{se $i \not\in S$}} \\
	1 &\text{\text{se $i \in S$}}
\end{cases}
	\]
\`E chiaro che:
\[ f(S) = f(T) \iff \forall {i \in \omega} \; a_i^S = a_i^T \iff \forall
{i \in \omega} \; i \in S \leftrightarrow i \in T \iff S = T
	\]
Non sarebbe difficile formalizzare questa dimostrazione.
Basterebbe definire $0.a_1a_2a_3\ldots$ come $\sum_{i = 0}^\infty a_i 10^{-(i + 1)}$, poi $\sum_{i = 0}^\infty$
come $\sup_{n \in \omega}\sum_{i = 0}^n$, poi $\sum_{i = 0}^n$ per ricorsione numerabile, poi dimostrare le propriet\`a aritmetiche rilevanti. Noi sfrutteremo la stessa idea, 
ma formulando la dimostrazione in termini di ordini.}

\subsection{Operazioni che coinvolgono la cardinalit\`a del continuo}
Prima di dimostrare il teorema, sviluppiamo un po' di aritmetica della
cardinalit\`a $2^{\aleph_0}$. {\drkred Questi lemmi sono importanti, e serviranno per calcolare 
la cardinalit\`a di insiemi concreti.}

\begin{remark}
	$(2^{\aleph_0})^{\aleph_0} = 2^{\aleph_0}$.
\end{remark}

\begin{proof}
$(2^{\aleph_0})^{\aleph_0} = 2^{\aleph_0 \cdot \aleph_0} = 2^{\aleph_0}$.
\end{proof}

\begin{lemma}
	\label{prodotto_cardinali_senza_AC}
	Siano $\alpha$, $\beta$ o ``finito'' o $\aleph_0$ o $2^{\aleph_0}$, allora:
	\[ \alpha + \beta = \alpha \cdot \beta = \max(\alpha,\beta)
		\]
	{\drkred eccetto il caso} $\alpha \cdot 0 = 0 \cdot \beta = 0$.
\end{lemma}

\begin{proof}
	Somme e prodotti di cardinalit\`a finite sono finiti
(teorema~\ref{op_card_fin}). 
Supponiamo quindi $\aleph_0 \leq \beta$ e, senza perdita di generalit\`a,
$\alpha \le \beta$. Allora
\[
\beta \;\;\le\;\; {\alpha+\beta \atop \alpha\cdot\beta} \;\;\le\;\;
\beta^2 \;\;{\drkred =}\;\;\beta
\]
	dove l'ultima uguaglianza vale perch\'e $\aleph_0^2 = \aleph_0$, e
$2^{\aleph_0} \leq (2^{\aleph_0})^2 \leq (2^{\aleph_0})^{\aleph_0}  =
2^{\aleph_0}$.
\end{proof}

\begin{lemma}
	Se $2 \leq \alpha \leq 2^{\aleph_0}$ allora $\alpha^{\aleph_0} = 2^{\aleph_0}$.
\end{lemma}

\begin{proof}
	$2^{\aleph_0} \leq \alpha^{\aleph_0} \leq (2^{\aleph_0})^{\aleph_0} = 2^{\aleph_0}
		$
\end{proof}

\subsection{$2^{\aleph_0} - \aleph_0 = 2^{\aleph_0}$}
Ricordiamo un'osservazione riguardo al numerabile.

\begin{remark}
	\label{numerabile - finito}
	Sia $|A| = \aleph_0$ e $B \subseteq A$ con $|B|<\aleph_0$. Allora $|A \setminus B| = \aleph_0$.
\end{remark}

\begin{proof}
	Siccome $A \setminus B \subseteq A$ o $|A \setminus B| = \aleph_0$ o $|A \setminus B|<\aleph_0$.
Se valesse la seconda $A = B \cup (A \setminus B)$ sarebbe
finito~$\;\lightning$
\end{proof}

Vale una proposizione analoga per $2^{\aleph_0}$.

\begin{lemma}
	\label{continuo-numerabile}
	Sia $|A| = 2^{\aleph_0}$ e $B \subseteq A$ con $|B| \leq \aleph_0$, allora $|A \setminus B| = 2^{\aleph_0}$.
\end{lemma}

\begin{note}
Assumendo l'assioma di scelta, il lemma vale anche rimpiazzando $|B| \leq
\aleph_0$ con $|B| < 2^{\aleph_0}$. Anzi, per qualunque~$A$ e
qualunque~$B\subseteq A$ purch\'e $|B|
< |A|$. Per\`o, per ora, possiamo dimostrare solo l'asserto pi\`u debole
enunciato qua sopra.
\end{note}

\begin{proof}
	Abbiamo la disuguaglianza $|A\setminus B| \leq |A| =
2^{\aleph_0}$, basta quindi stabilire che $2^{\aleph_0}\le |A\setminus B|$.
	Siccome $2^{\aleph_0}\cdot 2^{\aleph_0} = 2^{\aleph_0}$, esiste una bigezione:
	\[ f \colon A \to {}^{\omega}2 \times {}^{\omega}2
		\]
	sia $\pi \colon {}^{\omega}2 \times {}^{\omega}2 \to {}^{\omega}2$
la proiezione $\pi(x,y) = x$. Siccome $B$ \`e al pi\`u numerabile
	\[ |f[B]| \leq |B| \leq \aleph_0 < 2^{\aleph_0}
		\]
\vbox to 0pt{
\vskip-2.5\baselineskip
		\includegraphics[width = .3\textwidth]{immagini/continuo-numerabile.png}
\hss}

\vskip-\baselineskip
\vskip-\parskip
\setlength\leftskip{.33\textwidth}
	Per cui $\pi \circ f[B] \ne {}^{\omega}2$.
Prendiamo $x \in {}^{\omega}2 \setminus \pi \circ f[B]$.\\
Dire che $x \not \in \pi \circ f[B]$ significa
	che $(\{x\} \times {}^{\omega}2) \cap f[B] = \emptyset$.\\
	$\text{Quindi } f^{-1}(\{x\} \times {}^{\omega}2) \cap B = \emptyset$,\\
	$\mathrlap{\text{ossia}}\phantom{\text{Quindi }}f^{-1}(\{x\} \times {}^{\omega}2) \subseteq A \setminus B$,\\
	$\mathrlap{\text{da cui}}\phantom{\text{Quindi }}\mathllap{|}f^{-1}(\{x\} \times {}^{\omega}2)| \leq |A \setminus B|$.\\
	Usando il fatto che $f$ \`e bigettiva:\\
	$\phantom{\text{Quindi }}\mathllap{|}f^{-1}(\{x\} \times
{}^{\omega}2)| = |{}^{\omega}2| = 2^{\aleph_0}$.\\
Quindi $2^{\aleph_0}\le |A\setminus B|$.

\setlength\leftskip{0pt}
\vskip-\baselineskip
\end{proof}

\begin{proof}[Dimostrazione di~\ref{crdr} {\drkred(\mathshift|\RR| =
2^{\aleph_0}\mathshift)}]
	Siccome $\RR \subseteq \ps(\QQ)$, la disuguaglianza $|\RR| \leq
2^{\aleph_0}$ \`e immediata. Per dimostrare la disuguaglianza opposta
consideriamo:
	\[ A \;\;\Mydef\;\; \big\oldparen\;X \in \ps(\omega)\; |\; X \ne
\emptyset \land |\omega \setminus X| \geq \aleph_0\;\big\}
		\]
	{\drkgrn ossia i sottoinsiemi di $\omega$ non vuoti e \vocab{coinfiniti}.

	Intuitivamente: $X \in A$ rappresenta lo sviluppo in notazione
binaria di un $x \in\, ]0,1[$
\[ x = 0.a_1a_2a_3 \ldots \qquad \iff \qquad a_i = 1 \leftrightarrow i \in
X\]
La condizione $X \ne \emptyset$ serve a escludere lo $0$, la 
	condizione di coinfinitezza a escludere l'uno periodico.}

	Ci basta dimostrare che $|A| = 2^{\aleph_0}$ e che esiste $f
\colon A \rightarrow \RR$ iniettiva. La prima cosa \`e facile:
	\[ A = \ps(\omega)\;\setminus\; \Big(\;\{\emptyset\} \;\cup\;
\underbrace{\big\oldparen\: X \in \ps(\omega)\;\big|\; |\omega \setminus
X|<\aleph_0\:\big\}}_{=: S}\;\Big)
		\]
	L'insieme $S$ -- dei sottoinsiemi di $\omega$ \vocab{cofiniti} --
\`e in corrispondenza biunivoca con $\psf(\omega)$ tramite la funzione
\textit{complementare rispetto a $\omega$}:
	\[ \ol \bullet \colon S \rightarrow \psf(\omega) \qquad X \mapsto \ol X = \omega \setminus X
		\]
	Quindi $|S| = |\psf(\omega)| = \aleph_0$, e, di conseguenza $|A| = 2^{\aleph_0}$ grazie al lemma precedente.

	Resta da costruire $f \colon A \rightarrow \RR$ iniettiva.
	Cominciamo col dare un ordine totale su $A$. Dati $X,Y \in A$ definiamo:
	\[ X <_A Y \;\;\Mydef\;\; \exists {i \in \omega} \; \underbrace{(i
\cap X = i \cap Y)}_{\forall {j < i} \; j \in X \leftrightarrow j \in Y} \land \underbrace{(i \in Y \setminus X)}_{i \in Y \land i \not\in X}\
		\]
	ossia, detto $i$ il minimo della differenza simmetrica $X \Delta Y
\;\Mydef\; (X \setminus Y) \cup (Y \setminus X)$, allora per definizione o
$i\in Y$ o $i\in X$: nel primo caso diciamo che $X<Y$, nel secondo
che~$Y<X$.
\`E facile verificare che questo \`e un ordine totale.

Ora vorremmo immergere $A$ in $\RR$. L'idea \`e trovare un sottoinesieme
denso e numerabile di~$A$: questo sar\`a isomorfo a~$\QQ$, e la
completezza di~$\RR$ ci permetter\`a di estendere l'isomorfismo a una
iniezione di~$A$ in~$\RR$. Consideriamo $B \Mydef \psf(\omega) \setminus \{\emptyset\} \subseteq A$.
Sappiamo che $|B| = \aleph_0$. Dimostriamo che $B$ \`e denso in $A$.

\vbox to 0pt{
\vskip.5\baselineskip
		\includegraphics[width = 0.47\textwidth]{immagini/cardinalità_R.png}
\vss}

\vskip-\baselineskip
	Dati $X,Y \in A$, con $X < Y$, sia $i = \min X \Delta Y$ e sia
$j>i$
minimo tale che $j\not\in X$, che esiste perch\'e $X$ \`e coinfinito.
	Sia $Z = (X \cap j) \cup \{j\}$, che \`e finito, quindi appartiene
a~$B$. Per costruzione $j=\min X\Delta Z$ e $j\in Z$, quindi $X<Z$.
Inoltre $i=\min Y\Delta Z$ e $i\in Y$, quindi $Z<Y$. Abbiamo cos\`i
verificato la densit\`a.
\setlength\leftskip{.5\textwidth}

\setlength\leftskip{0pt}
	Stabilito che $B$ \`e denso in $A$,  allora, in particolare, $B$ \`e denso in se stesso, dunque
 	valgono le ipotesi del \hyperref[isoCantor]{teorema di isomorfismo
di Cantor}, quindi c'\`e $g \colon B \rightarrow \QQ$, isomorfismo di ordini.
	Ora, siccome $B$ \`e denso in $A$, la funzione:
	\[ h \colon A \rightarrow \; \{\text{sezioni su $B$}\} \qquad X \mapsto B_X = \{Y \in B | Y < X\}
		\]
	\`e ben definita e iniettiva. Quindi $f \colon A \rightarrow \RR$
definita da~$f(X) = g[h(X)]$ \`e iniettiva.
\end{proof}
%
%\subsection{\texorpdfstring{$(\star)$ Alternativa per la cardinalit\`a di $\RR$}{Alternativa per la cardinalit\`a dei reali}}
%Vediamo un modo alternativo di stimare dal basso la cardinalit\`a di $\RR$ sfruttando l'insieme di Cantor.
%L'idea \`e che ad successione binaria possiamo associare un punto dell'insieme di Cantor scegliendo semplicemente il percorso sui rami dell'albero della costruzione,
%dopodich\'e - poich\'e intersechiamo una successione decrescente di compatti - otteniamo un insieme non vuoto da cui prendere il punto di $C$ da far corrispondere alla sequenza binaria.
%\begin{figure}[H]
%	\centering
%	\includegraphics[scale = 0.4]{immagini/cantor-R-cardinalità.png}
%\end{figure}
%In tal modo avremo $2^{\aleph_0} = |2^{\omega}| \leq |C| \leq |\RR|$.
%
%\begin{proof}
%	Definiamo la seguente funzione:
%	\[ f : {}^{\omega}2 \to \RR : (a_i)_{i \in \omega} \mapsto \sum_{i = 0}^{\infty} \frac{a_i}{3^i} \;\textcolor{MidnightBlue}{= a_1 \cdot 3^{-1} + a_2 \cdot 3^{-2} + \ldots}
%		\]
%	\textcolor{MidnightBlue}{- sarebbe una scrittura in base 3 senza la cifra 2\footnote{Potremmo anche fare la stessa cosa buttando
%	via 1 anzich\'e 2, ma convenzionalmente si sceglie che \`e meglio avere 0.1 che 0.0222\ldots. Potremmo altres\`i usare una base diversa ottenendo la stessa cosa, purch\'e non sia due, in tal caso si procede come la prima dimostrazione vista.}, cio\`e una cosa del tipo $0.01011100\ldots$,
%	dove non abbiamo problemi di ambiguit\`a a causa di periodicit\`a, dati da $0,1 = 0,02222\ldots$, poich\'e abbiamo tolto il 2 e quindi usiamo solo la prima scrittura. In particolare \`e come elencare i punti del Cantor fatto in $\left[0,\frac{3}{2}\right]$.}\footnote{Infatti si poteva procedere anche usando il Cantor solito e definendo la mappa da $\{0,2\}^\omega$ a $\RR$,
%	che manda una successione in $\sum_{i = 1}^\infty \frac{a_i}{3^i} = \frac 13 \sum_{i = 0}^\infty \frac{a_i}{3^i}$, e in particolare, questa si ottiene dalla funzione precedente moltiplicata per $\frac{2}{3}$ (notare che questa cosa cambia l'insieme numerico da $\{0,1\}$ a $\{0,2\}$ e sistema l'indice per il primo termine della scrittura decimale).}
%	\begin{figure}[H]
%		\centering
%		\includegraphics[scale = 0.4]{immagini/cantor_0_32.png}
%	\end{figure}
%	Dove naturalmente la sommatoria \`e definita per ricorsione numerabile come funzione da $\omega \to \RR$ nella maniera seguente:
%	\[ \sum_{i = 0}^0 \square_i = \square_0 \qquad \sum_{i = 0}^{n + 1} \square_i = \left(\sum_{i = 0}^n \square_i\right) + \square_{n + 1}
%		\]
%	e la serie \`e ben definita come:
%	\[ \sum_{i = 0}^{\infty}\square_i = \sup_{n \in \omega}\left(\sum_{i = 0}^n \square_i\right)
%		\]
%	Occorre quindi verificare che $f$ \`e ben definita ed iniettiva.
%	\begin{itemize}
%		\item[$\diamondsuit$] \underline{$f$ ben definita}: cio\`e la serie in arrivo \`e un elemento di $\RR$, ossia $\sup_{n \in \omega}\left(\sum_{i = 0}^n \frac{a_i}{3}\right)$ esiste,
%		e, come sappiamo, ci\`o \`e equivalente a chiedere che $\left\{\sum_{i = 0}^n \frac{a_i}{3}\right\}_{n \in \omega}$ ha un maggiorante. Vediamo per induzione che:
%		\[ \sum_{i = 0}^n \frac{a_i}{3^i} \leq \frac 32 - \frac{1}{2 \cdot 3^n}
%			\]
%		fatto questo il RHS si pu\`o stimare a sua volta con $\frac 32$ \textcolor{MidnightBlue}{- sarebbe la somma della serie -}, che costituisce un upper bound per ogni $n \in \omega$, quindi per tutti gli elementi dell'insieme.
%		\begin{itemize}
%			\item[$\boxed{n = 0}$] Banalmente $ \sum_{i = 0}^0 \frac{a_i}{3^i} = \frac{a_0}{3^0} = \frac{a_1}{3} \leq \frac 13 < \frac 43 = \frac 32 - \frac{1}{2 \cdot 3^1}$.
%			\item[$\boxed{n \implies n + 1}$] Supponiamo $\sum_{i = 0}^n \frac{a_i}{3^i} \leq \frac 32 - \frac{1}{2 \cdot 3^n}$ e dimostriamo $\sum_{i = 0}^{n+1} \frac{a_i}{3^i} \leq \frac 32 - \frac{1}{2 \cdot 3^{n+1}}$.
%			\begin{align*}
%				\sum_{i = 0}^{n+1} \frac{a_i}{3^i} &= \left(\sum_{i = 0}^{n} \frac{a_i}{3^i}\right) + \frac{a_{n+1}}{3^{n+1}} &&\text{(def. ricorsiva)} \\
%												   &\leq \frac 32 - \frac{1}{2 \cdot 3^n} + \frac{a_{n+1}}{3^{n+1}} &&\text{(hp. induttiva)} \\
%												   &\leq \frac 32 - \frac{1}{2 \cdot 3^n} + \frac{1}{3^{n+1}} &&(a_{n+1} \in \{0,1\}) \\
%												   &= \frac 32 - \frac{1}{3^n} \cdot \frac 16 = \frac 32 - \frac{1}{2 \cdot 3^{n+1}}
%			\end{align*}
%		\end{itemize}
%		\item[$\diamondsuit$] \underline{$f$ iniettiva}: date due successioni binarie distinte, $(a_i)_{i \in \omega}$ e $(b_i)_{i \in \omega}$, sia $i_0 \in \omega$ il minimo per cui $a_{i_0} \ne b_{i_0}$. Assumiamo WLOG di essere nel caso $a_{i_0} = 0$ e $b_{i_0} = 1$ e 
%		verifichiamo che $f(b) > f(a)$ \textcolor{MidnightBlue}{- stiamo proprio verificando la stretta crescenza}. In particolare dimostriamo che:
%		\[ f(a) < f(a) + \frac{1}{2 \cdot 3^{i_0}} \leq f(b)
%			\]
%		Per verificare la seconda disuguaglianza osserviamo che:
%		\begin{align*}
%			f(a) &= \sup_{n \in \omega} \left(\sum_{i = 0}^n \frac{a_i}{3^i}\right) \\
%				 &= \sup_{n \in \omega} \left(\sum_{i < i_0}\frac{a_i}{3^i} + \sum_{i = i_0 + 1}^n \frac{a_i}{3^i} \right) &&(a_{i_0 + 1} = 0) \\
%				 &= \sup_{n \in \omega} \left(\sum_{i < i_0}\frac{b_i}{3^i} + \sum_{i = i_0 + 1}^n \frac{a_i}{3^i} \right) &&(\forall i < i_0 \; a_{i} = b_{i}) \\
%				 &=\sup_{n \in \omega} \left(\sum_{i < i_0}\frac{b_i}{3^i} + \frac{1}{3^{i_0 + 1}}\sum_{j = 0}^{n - i_0 - 1} \frac{a_{j + i_0 + 1}}{3^j} \right) &&(j = i - i_0 - 1) \\
%				 &\leq\sup_{n \in \omega} \left(\sum_{i < i_0}\frac{b_i}{3^i} + \frac{1}{3^{i_0 + 1}} \cdot \frac 32  \right) &&(\text{dalla stima sopra}) \\
%				 &= \sum_{i < i_0}\frac{b_i}{3^i} + \frac{1}{2 \cdot 3^{i_0}}  
%		\end{align*}
%		Ora, dato che:
%		\begin{align*}
%			f(b) &\geq \sum_{i < i_0}\frac{b_i}{3^i} + \frac{1}{3^{i_0}} \implies \sum_{i < i_0}\frac{b_i}{3^i} \leq f(b) - \frac{1}{3^{i_0}} &&(b_{i_0} = 1)
%		\end{align*}
%		possiamo concludere la stima sopra:
%		\[ f(a) \leq \sum_{i < i_0}\frac{b_i}{3^i} + \frac{1}{2 \cdot 3^{i_0}} \leq f(b) - \frac{1}{3^{i_0}} + \frac{1}{2 \cdot 3^{i_0}} = f(b) - \frac{1}{2 \cdot 3^{i_0}}
%			\]
%		che \`e equivalentemente alla seconda disuguaglianza voluta.
%	\end{itemize}
%\end{proof}
%
%\subsection{\texorpdfstring{$(\star) \; (F,0,1,+,\cdot,\leq)$}{Unicit\`a (a meno di isomorfismo) di un campo ordinato completo}}
